% Speaker script: slide-by-slide talking points for the 23 Nov presentation, plus Q&A prep.
% Companion to slides.tex (same content as its \note{} text, broken into short lines to present from).
% Build: `make speaker_script.pdf` (or latexmk -pdf speaker_script.tex).
\documentclass[12pt]{article}
\usepackage[letterpaper,margin=0.8in,top=0.7in]{geometry}
\usepackage[T1]{fontenc}
\usepackage[sfdefault,lining]{FiraSans}
\usepackage{newtxsf}
\usepackage{xcolor}
\usepackage{enumitem}
\usepackage[most]{tcolorbox}
\usepackage{titlesec}
\usepackage{hyperref}
\hypersetup{pdftitle={Speaker script: Prediction Horizon and the Content of Neural Population Subspaces},
  pdfauthor={Gunnar Andra-Thomas}, colorlinks=true, linkcolor=black}

\definecolor{teal}{HTML}{23373B}
\definecolor{accent}{HTML}{2A78D6}
\definecolor{ink}{HTML}{52514E}
\definecolor{warn}{HTML}{B4521E}

\setlength{\parindent}{0pt}
\setlength{\parskip}{4pt}
\setlist[itemize]{leftmargin=1.3em, itemsep=2pt, topsep=2pt, parsep=0pt}
\setlist[itemize,2]{leftmargin=1.2em, label=\textendash}
\titleformat{\section}{\Large\bfseries\color{teal}}{}{0pt}{}
\titlespacing*{\section}{0pt}{0pt}{6pt}
\titleformat{\subsection}{\large\bfseries\color{teal}}{}{0pt}{}
\titlespacing*{\subsection}{0pt}{10pt}{3pt}

% \slide{number}{title}{time on slide}{clock at end of slide}
\newcommand{\slide}[4]{%
  \newpage
  \begin{tcolorbox}[colback=teal, colframe=teal, coltext=white, arc=2pt, boxsep=3pt, left=6pt, right=6pt]
    {\small SLIDE #1 \hfill #3 on this slide \textperiodcentered{} finish by #4}\\[2pt]
    {\Large\bfseries #2}
  \end{tcolorbox}\vspace{2pt}}
\newcommand{\point}[1]{{\color{ink}\textbf{On screen:} #1}\par}
\newcommand{\say}{\par\smallskip{\color{accent}\textbf{Say}}\par}
\newcommand{\avoid}[1]{\par\smallskip{\color{warn}\textbf{Careful:} #1}\par}
\newcommand{\bridge}[1]{\par\smallskip\begin{tcolorbox}[colback=accent!8, colframe=accent!40, arc=2pt, boxsep=2pt,
  left=5pt]\textbf{Bridge to next slide:} #1\end{tcolorbox}}
\newcommand{\ifasked}[1]{\par\smallskip{\color{ink}\textbf{Only if asked:} #1}\par}

% Q&A entries: \q{question}{bullets}
\newcommand{\q}[2]{\par\medskip{\bfseries\color{teal}Q: #1}\par\vspace{1pt}\begin{itemize}#2\end{itemize}}

\begin{document}

% ======================================================================================
\begin{center}
{\LARGE\bfseries\color{teal} Speaker script}\\[4pt]
{\large Prediction Horizon and the Content of Neural Population Subspaces}\\[2pt]
{\color{ink}Gunnar Andra-Thomas \textperiodcentered{} CSE 8803 NBD \textperiodcentered{} Monday 23 November 2026}
\end{center}

\subsection*{How to use this}
\begin{itemize}
\item One page per slide. The box at the top gives the slide number, title, time budget and the clock time you
should reach by the end of that slide.
\item \textbf{Say} lines are talking points, not a script to read. Say each one in your own words.
\item Numbers on the slides are there for questions. Say only the few marked here.
\item 12 content slides plus the title, about 15 minutes. If you're behind at slide 8, shorten slides 9 and 10.
\end{itemize}

\subsection*{The whole talk in four sentences}
\begin{enumerate}[leftmargin=1.5em, itemsep=2pt]
\item I asked whether predicting brain activity further ahead pulls out the animal's slow belief about the task.
\item It doesn't: there is no detectable rise in block-prior information with prediction horizon.
\item What every low-dimensional summary picks up instead is a slow change across the session.
\item These methods rank directions by how much variance they predict, and that big slow signal wins first.
\end{enumerate}

\subsection*{Three distinctions to keep straight all talk}
\begin{itemize}
\item \textbf{Corrected score} = real score minus the pseudo-session null median. Zero means ``no advantage over the
null''. It is \emph{not} percent of variance explained.
\item \textbf{Prior:} ``no detectable rise''. Never say ``it falls''.
\item \textbf{Stimulus:} long horizons \emph{lose} it. Short horizons are not \emph{specialized} for it.
\end{itemize}

\subsection*{Clock checkpoints}
\begin{itemize}
\item End of slide 4 (method): about 4:05.
\item End of slide 8 (prior): about 9:35.
\item End of slide 13 (takeaways): about 15:00.
\end{itemize}

% ======================================================================================
\slide{1}{Title}{20 s}{0:20}
\point{Title, your name, course.}
\say
\begin{itemize}
\item Hi, I'm Gunnar.
\item The question: when you look for the directions of brain activity that best predict its own future, does
\emph{how far ahead} you predict decide what they encode?
\item Short answer, in mouse secondary motor cortex: mostly not.
\item Data: public Neuropixels recordings from the International Brain Laboratory.
\item The three main tests were pre-registered before I fitted real data.
\end{itemize}

% ======================================================================================
\slide{2}{Is a slow belief among the dominant predictable signals?}{1:00}{1:20}
\point{Two columns: Neuroscience (left) and Methods (right).}
\say
\begin{itemize}
\item \textbf{Left, neuroscience.}
  \begin{itemize}
  \item Animals track hidden task states over many trials.
  \item The IBL ``block prior'' is a well-studied example: it can be decoded weakly across the brain (Findling 2025).
  \item Question: is a belief like this one of the \emph{dominant} predictable signals in population activity?
  \end{itemize}
\item \textbf{Right, methods.}
  \begin{itemize}
  \item Predictive methods find structure without labels: lagged regression, slow features, DCA, and predictive
  learning in ML.
  \item If the prediction horizon sorted signals by timescale, long horizons would find belief states for free.
  \end{itemize}
\item Preview of the punchline: these methods care about more than timescale. They rank directions by how much
future variance they explain.
\end{itemize}

% ======================================================================================
\slide{3}{Hypothesis: the prediction horizon sorts content by timescale}{1:15}{2:35}
\point{Grey step line = the true block. Blue curve = the ideal-observer prior. Hypothesis box on the right.}
\say
\begin{itemize}
\item \textbf{The task.} A grating appears left or right; the mouse turns a wheel to report the side.
\item The side is biased 80/20, in blocks that switch \emph{without a cue}. So the mouse has to infer the block.
\item \textbf{Grey line:} the true block.
\item \textbf{Blue curve:} the Bayes-optimal prior, what an ideal observer would believe. It is a slow variable
that lags the true block.
  \begin{itemize}
  \item It is a model of the belief, not a measurement from the brain.
  \end{itemize}
\item \textbf{Hypothesis:}
  \begin{itemize}
  \item Predicting 30 s ahead isolates a slow subspace carrying the prior.
  \item Predicting 0.5 s ahead isolates a fast subspace carrying the current stimulus.
  \end{itemize}
\item \textbf{Three tests, fixed before any real data:}
  \begin{itemize}
  \item Geometry: the subspaces differ.
  \item Prior: decoding rises with horizon.
  \item Stimulus: decoding falls with horizon.
  \end{itemize}
\item The prior and stimulus tests are the heart of it. Geometry only asks whether the subspaces differ at all.
\end{itemize}

% ======================================================================================
\slide{4}{Method: predict activity $\Delta$ s ahead through a 5-dimensional bottleneck}{1:30}{4:05}
\point{Diagram: activity now $\to$ 5 latents $\to$ activity $\Delta$ later. The formula underneath.}
\say
\begin{itemize}
\item \textbf{Data line:} mouse secondary motor cortex (MOs), 16 sessions from 13 mice, activity in half-second bins.
\item \textbf{The picture.}
  \begin{itemize}
  \item Take the population now and squeeze it through 5 latents.
  \item Use those to predict the population $\Delta$ seconds later.
  \item Four horizons: 0.5, 2, 8 and 30 s.
  \end{itemize}
\item \textbf{The formula says the same thing.} Point to each part as you say it:
  \begin{itemize}
  \item First term: how badly activity now predicts activity $\Delta$ later.
  \item The rank limit of 5: the bottleneck.
  \item $\lambda$: a ridge penalty, chosen by cross-validation once per session and shared by all lags.
  \end{itemize}
\item Name: lagged reduced-rank regression. PCA is the unsupervised baseline.
\item \textbf{What gets fitted:} the 5 directions of current activity that best predict the future. That is
``the subspace''.
\item \textbf{What gets decoded:} from those 5 latents, I decode the stimulus, the prior and elapsed time with simple
cross-validated decoders.
\item Null and statistics are in the backup slides.
\end{itemize}
\bridge{``Before testing the hypothesis, two sanity checks.''}

% ======================================================================================
\slide{5}{The prior is in MOs; 5 dimensions capture most predictable activity}{1:15}{5:20}
\point{Left: three big numbers (prior from the full population). Right: line plot of prediction accuracy vs.\ lag.}
\say
\begin{itemize}
\item \textbf{Check 1, left.}
  \begin{itemize}
  \item Decoding the prior from the whole population, the way Findling did, beats the pseudo-session null.
  \item But it is small: about 0.06.
  \item Keep that number in mind; it matters later.
  \end{itemize}
\item \textbf{Check 2, right.}
  \begin{itemize}
  \item Each line is a method. The vertical axis is how well it predicts held-out future activity.
  \item Every line falls further ahead, as expected.
  \item Lagged RRR (blue) beats PCA and random subspaces of the same size.
  \item At 30 s, 5 dimensions get 96\% of an unconstrained model.
  \end{itemize}
\item \textbf{The caveat. Say this one carefully:}
  \begin{itemize}
  \item 5 dimensions capture nearly all the \emph{aggregate} predictability.
  \item That does not guarantee they capture small task signals. A signal that explains little predictable variance
  can be left out.
  \end{itemize}
\end{itemize}
\bridge{``So the method works. Now the three tests, starting with geometry.''}

% ======================================================================================
\slide{6}{Subspaces differ by lag, but 30 s ones don't reproduce across halves}{1:30}{6:50}
\point{Three bars (same lag 0.5 s, cross lag, same lag 30 s) and the dashed chance line at about 68\textdegree.}
\say
\begin{itemize}
\item \textbf{How to read it.}
  \begin{itemize}
  \item Angle between subspaces: 0\textdegree{} = they share a direction; 90\textdegree{} = completely different.
  \item I fit each lag separately on the two halves of a session.
  \item Each bar is the angle between those fits.
  \item Dashed line: chance, i.e.\ random directions with the same variance profile. It is already about
  68\textdegree.
  \end{itemize}
\item \textbf{Tell it as a story. Intuition first, numbers second:}
  \begin{itemize}
  \item We want to know whether short- and long-lag subspaces differ.
  \item First, how reproducible is each one?
  \item Short-lag fits reproduce: left bar, about 59\textdegree, well below chance.
  \item 30 s fits do not: right bar, at chance, even for their single closest direction.
  \item The cross-lag angle beats both, so the pre-registered test passes ($p = 0.0026$).
  \end{itemize}
\item \textbf{But it passes mostly through the short lag.} So it is \emph{not} two clean, stable representations.
\item This is the weakest of the results.
\end{itemize}
\avoid{don't read the angles one by one. ``Short reproduces, long is at chance'' is the message.}
\bridge{``That's whether they differ. Now what they contain, starting with the stimulus.''}

% ======================================================================================
\slide{7}{Longer horizons lose the current stimulus}{1:15}{8:05}
\point{Bar chart: blue = RRR at 0.5, 2, 8, 30 s; grey = PCA and random.}
\say
\begin{itemize}
\item I show the stimulus first because it is the cleaner positive result.
\item \textbf{How to read it.}
  \begin{itemize}
  \item Each bar: how well a subspace decodes the current stimulus side.
  \item Minus what the same decoder gets on pseudo-sessions.
  \item So zero = no advantage over the pseudo-session null.
  \end{itemize}
\item \textbf{Result.}
  \begin{itemize}
  \item Decoding falls as the horizon grows ($p = 0.006$, 10 of 13 mice).
  \item It holds after controlling for block information.
  \end{itemize}
\item \textbf{The nuance:} the 0.5 s subspace is no better than a random one (grey).
  \begin{itemize}
  \item So long lags \emph{lose} the stimulus.
  \item Short lags don't \emph{specialize} for it.
  \end{itemize}
\item The stimulus half of the hypothesis holds, in a weaker form than predicted.
\item Holds at 10 dimensions, not at 3: real, but modest.
\end{itemize}
\bridge{``Now the central prediction: the prior.''}

% ======================================================================================
\slide{8}{Predicting further ahead: no detectable gain in block-prior information}{1:30}{9:35}
\point{Left: same bar layout for the prior. Right: one dot per mouse = its change from 0.5 s to 30 s.}
\say
\begin{itemize}
\item Same layout. The hypothesis predicted the blue bars to \emph{rise} left to right.
\item \textbf{Right panel.}
  \begin{itemize}
  \item Each dot is one mouse's change from 0.5 to 30 s.
  \item Already corrected against that mouse's own pseudo-sessions.
  \item Most sit at or below zero. Only 4 of 13 rise.
  \end{itemize}
\item \textbf{The left bars} at 2--30 s have intervals that include zero: noise around the null.
\item \textbf{This was the central prediction, and it is not supported.}
\item The prior \emph{is} real in the population (slide 5). It is just a small signal.
\item \textbf{The wording to use:} ``no detectable rise''.
  \begin{itemize}
  \item Individual bars move by about $\pm0.05$ if you change the decoder or detrend.
  \item The one thing that holds in every version is that there's no rise.
  \end{itemize}
\end{itemize}
\avoid{don't say the prior ``falls'' or that the 0.5 s subspace ``has the prior''. The 0.054 is distance above the null,
not 5\% of variance.}
\bridge{``So if not the prior, what does the 30 s subspace carry?''}

% ======================================================================================
\slide{9}{The leading 30 s direction tracks session time, and so does PCA}{1:15}{10:50}
\point{4$\times$4 grid: all 16 sessions. Blue = activity along the leading 30 s direction; dashed = straight-line fit.}
\say
\begin{itemize}
\item \textbf{Each small panel is one session.}
  \begin{itemize}
  \item Blue: activity along the subspace's leading direction, in 30 s windows, across the session.
  \item Panels are sorted by how closely it follows time.
  \end{itemize}
\item In 12 of 16 it is nearly a straight ramp. A few look more like steps.
\item Showing all 16 means you see the typical case, not a picked example.
\item \textbf{Elapsed-time decoding} (bottom line):
  \begin{itemize}
  \item 30 s subspace: 0.84.
  \item But the 0.5 s subspace, PCA, and even random subspaces are close behind.
  \end{itemize}
\item \textbf{Message:} every low-dimensional summary of these recordings is dominated by this session ramp.
Long horizons have it only slightly more.
\end{itemize}
\bridge{``The obvious worry is that this is recording drift.''}

% ======================================================================================
\slide{10}{Recording controls leave most of the ramp, except the strictest unit filter}{1:15}{12:05}
\point{Each row = one version of the data, refit from scratch. Filled dot = control; hollow = original on the same
sessions; dashed lines = pre-registered 75\% and 50\% thresholds.}
\say
\begin{itemize}
\item Each row: I changed the data to remove one possible cause, then refit everything.
\item \textbf{These each leave most of the ramp:}
  \begin{itemize}
  \item removing trial-locked responses;
  \item dropping units that fade in and out;
  \item skipping spike sorting entirely (multi-unit activity);
  \item correcting for spikes shrinking in amplitude.
  \end{itemize}
\item Spike-amplitude coupling actually has the \emph{wrong sign} for an amplitude artifact.
\item Probe motion, pupil, engagement and satiety: no link beyond the shared trend (backup slides).
\item \textbf{The exception, and the biggest open challenge:} the strictest filter.
  \begin{itemize}
  \item Units stable in both rate and amplitude: only 5 sessions pass.
  \item The ramp falls by about a third, and PCA falls just as much.
  \item Is that just fewer units, or removing the units that carry it? Matched-unit-count resampling would tell.
  \end{itemize}
\end{itemize}
\bridge{``Is this specific to MOs?''}

% ======================================================================================
\slide{11}{The same pattern holds in other areas on the same probes}{1:00}{13:05}
\point{Left: elapsed-time bars, MOs (blue) vs.\ other areas (orange). Right: the three tests.}
\say
\begin{itemize}
\item The probes pass through other areas too.
\item In each session I took the largest other area (mostly frontal) and reran the \emph{unchanged} pipeline.
\item Criteria were fixed before binning.
\item \textbf{Left:} same ordering, slightly lower. The ramp is everywhere.
\item \textbf{Right:}
  \begin{itemize}
  \item Geometry and stimulus effects hold there too.
  \item The prior shows no rise, but it is barely present in these areas, so that part says little.
  \end{itemize}
\item \textbf{Honest limit:} same probes, so this is corroboration, not independent replication.
\end{itemize}

% ======================================================================================
\slide{12}{What the data show, possible mechanisms, and the main limits}{1:10}{14:15}
\point{Left = what the data show. Right = speculation. Bottom line = limits.}
\say
\begin{itemize}
\item \textbf{Left, the data:}
  \begin{itemize}
  \item no detectable rise of the prior with horizon;
  \item one session-time direction dominates, shared by PCA and short-lag RRR;
  \item no single recording cause tested explains most of it.
  \end{itemize}
\item \textbf{Right, speculation. Most plausible: variance ranking.}
  \begin{itemize}
  \item These methods keep the directions that predict the most future variance.
  \item A 30 s lag spans only a few trials, and blocks last minutes.
  \item So a persistent prior \emph{would} be predictable at every horizon tested.
  \item But it is small, and with 5 dimensions the big ramp wins first.
  \item Plausible, not demonstrated.
  \end{itemize}
\item \textbf{Limits:}
  \begin{itemize}
  \item the strictest unit filter;
  \item 13 mice, and the other areas share probes;
  \item the follow-ups are exploratory.
  \end{itemize}
\end{itemize}

% ======================================================================================
\slide{13}{Takeaways}{45 s}{15:00}
\point{Three numbered takeaways and a ``Next'' line.}
\say
\begin{itemize}
\item \textbf{One.} The predicted enrichment of the block prior at longer horizons was not supported.
\item \textbf{Two.} The dominant predictable direction tracks session time, and it is not special to long horizons.
\item \textbf{Three} (the one to remember). Its origin is unresolved despite targeted controls. Detrending did not
produce a robust gain in block-prior information at longer horizons.
\item \textbf{Lesson, within this dataset:} prediction horizon alone did not preferentially recover the block prior.
\item \textbf{Next:}
  \begin{itemize}
  \item matched-unit controls;
  \item a ground-truth simulation, to check the method can find a small slow signal;
  \item independent recordings.
  \end{itemize}
\item ``Thanks, happy to take questions.''
\end{itemize}

% ======================================================================================
\newpage
\section*{Backup slides: what each one is for}
Jump to these only when a question calls for them.
\begin{itemize}[itemsep=4pt]
\item \textbf{Data, null, statistics, pre-registration.} Region choice, session counts, how the pseudo-session null
works, Wilcoxon and Holm, which analyses were pre-registered.
\item \textbf{The factorization, and four things to explain.} The RRR closed form, readout vs.\ loading basis,
pseudo-sessions, the intersection--union test.
\item \textbf{Why lagged RRR rather than PCA or other methods.} CCA, DCA, GPFA, slow features; RRR's modest edge.
\item \textbf{How strong is the prior result?} Every variant of the prior test; the ``no detectable rise'' wording.
\item \textbf{Robustness table.} Holm $p$ for all three tests across 10 variants.
\item \textbf{Full decoding matrix.} Every variable by every subspace.
\item \textbf{Nesting.} The 0.5 s subspace contains 92\% of the 30 s leading direction.
\item \textbf{Candidate recording explanations.} One row per cause tested.
\item \textbf{Amplitude pilot.} Wrong-sign coupling; corrected refits keep 93--95\%.
\item \textbf{Probe motion, sorting, state signals.} Motion, pupil, engagement, satiety.
\item \textbf{The normalization bug.} The one post-results change.
\end{itemize}

% ======================================================================================
\newpage
\section*{Q\&A prep}
\textbf{Answer pattern:} one-line answer first, then the reason, then stop. If you don't know: ``I didn't test that;
what would settle it is \ldots''

\subsection*{The hypothesis and the setup}

\q{Why would horizon sort content by timescale?}{
\item A signal can only predict the future if it's still there.
\item A stimulus response lasts under a second; a block belief lasts minutes.
\item The catch: long horizons keep \emph{all} slow signals, ranked by size, and the biggest one here isn't the
belief.}

\q{Why MOs?}{
\item A rule fixed before any analysis chose it: at least 15 sessions from at least 8 mice.
\item MOs was the only area that qualified (16 sessions, 13 mice; MOp had 14 from 12).
\item It's also a frontal area associated with expectation signals.}

\q{The prior is a model. How do you know the mouse computes it?}{
\item I don't claim it does. It's a reference variable.
\item The question is whether activity tracks it beyond what slow drift alone would give. The pseudo-session null tests
exactly that.
\item The mice's choices do shift with the block.}

\q{Why such long horizons? Why not minutes, to match the blocks?}{
\item Sessions are about 66 minutes, so a 30 s lag already leaves only about 130 independent windows.
\item And blocks last minutes, so the prior is already predictable at 30 s. Horizon length wasn't the limit.}

\subsection*{Method}

\q{Why RRR instead of PCA?}{
\item It makes the horizon an explicit knob: only $\Delta$ changes between fits.
\item Closed form and well studied.
\item Honestly, it beats PCA only modestly, and much of that edge is trial-locked structure.}

\q{Why not DCA, slow feature analysis, CCA or GPFA?}{
\item DCA predicts over a window, not at a chosen horizon. Slow features and CCA are close cousins. GPFA was a planned
extension.
\item I'd expect any predictability-ranked method to find the ramp first. Comparing them is a next step.}

\q{Why 5 dimensions?}{
\item Fixed in advance.
\item 3 and 10 give the same pattern, except the stimulus test fails at 3.
\item At 30 s, 5 dimensions already get 96\% of full predictability.}

\q{Why one ridge penalty shared across lags?}{
\item So any difference between lags comes from the lag, not the penalty.
\item Still chosen by cross-validation.}

\q{What's the readout basis vs.\ the loading basis?}{
\item Readout: which current activity the model reads from (the registered subspace).
\item Loading: which future patterns it predicts.
\item Geometry passes in readout and fails in loading ($p = 0.96$), consistent with one shared slow signal.}

\q{Why contiguous cross-validation folds?}{
\item Neighbouring bins are highly correlated, so random folds leak the future into training.
\item Contiguous, purged folds are stricter. That's also why raw prior scores go negative.}

\subsection*{Null and statistics}

\q{What's a pseudo-session, and why is it a valid null?}{
\item A fake block sequence drawn from the experimenter's own rules.
\item Slow neural drift correlates with fake blocks just as well as real ones, but the brain can't know the fake ones.
\item Beating them means more than chance alignment with drift.}

\q{Why not just shuffle trials?}{
\item Shuffling destroys the slow structure, so the null becomes too easy and any drift looks like ``prior coding''.}

\q{Does a corrected 0.054 mean 5\% of variance?}{
\item No. It's the distance above the null median on that decoder's scale.
\item Zero means no advantage over the null.}

\q{Why are the raw prior R-squared values negative?}{
\item Contiguous folds: a decoder trained on one stretch extrapolates badly to another.
\item The null suffers equally, so the corrected difference is still meaningful.}

\q{Why the mouse as the unit? Why one-sided tests?}{
\item Sessions from one mouse aren't independent, so they're averaged first.
\item One-sided because the hypotheses are directional. That was fixed before the data.
\item Two-sided values are still significant: 0.0051 and 0.012.}

\q{Enough power with 13 mice?}{
\item For moderate effects, yes: geometry and stimulus are under 0.01.
\item For a small prior rise, no. That's why I say ``no detectable rise'', not ``no rise''.}

\subsection*{Results}

\q{Isn't ``no rise'' just absence of evidence?}{
\item Yes, and that's how I phrase it.
\item It's consistent: only 4 of 13 mice rise, and no rise shows up in any variant, decoder or detrend.
\item A modest rise isn't excluded: after the cosine detrend the interval reaches +0.068.}

\q{The 0.5 s bar (0.054) is higher than 30 s. So the prior falls?}{
\item I don't claim that. The single bars are at the noise floor and move about $\pm0.05$ with the decoder.
\item With Findling's decoder, 30 s is slightly \emph{higher} (0.038 vs 0.026).}

\q{How much of the prior do the subspaces capture?}{
\item It depends on the decoder.
\item Same decoder on the full population gives 0.159: 0.5 s captures about a third, PCA 8\%, 30 s none.
\item With Findling's decoder, 30 s captures about half.
\item Only the horizon comparison is robust.}

\q{Why is the 0.5 s subspace no better than random for the stimulus?}{
\item Stimulus information is spread widely, so almost any 5 directions pick some up.
\item The effect is ``long lags lose it'', not ``short lags specialize''.}

\q{If one slow mode dominates, why don't the 30 s halves agree?}{
\item The time course is consistent, but which units carry it changes within a session.
\item Post hoc: the leading direction isn't reproducible in filters, loadings, activation patterns or capture.
\item Caveat: the halves are different times, so reproducibility and change over time are confounded.}

\q{How does a random subspace decode elapsed time at 0.64?}{
\item The ramp is large and spread across many units, so almost any direction picks it up.
\item That's the point: session time dominates every low-dimensional view.}

\q{Elapsed-time decoding is just interpolation within a session, isn't it?}{
\item Yes. It measures how smoothly a representation varies with time.
\item All subspaces are scored the same way. The comparison is what matters.}

\subsection*{The session-time mode}

\q{Isn't it just electrode drift?}{
\item I tested the main candidates one by one: unit turnover, spike sorting, rigid probe motion, spike amplitude.
\item None explains most of it.
\item The strictest unit filter weakens it (0.73 to 0.49 in 5 sessions; PCA falls equally). That's the open
challenge.
\item So: not established as neural, not explained by any single artifact I tested.}

\q{Why is the amplitude coupling negative?}{
\item During fast firing and bursts, spikes get smaller: rate affects amplitude.
\item Detection loss would give the opposite sign.
\item Still negative using isolated spikes only.}

\q{Why not just remove the drift and refit?}{
\item I did, post hoc, two ways.
\item Slow-cosine detrend: a small prior appears about equally in every subspace, PCA included.
\item Quadratic detrend: the prior is near zero everywhere, and the stimulus effect weakens.
\item Neither shows a detectable rise with horizon.}

\q{Other areas on the same probe: isn't that the same artifact twice?}{
\item It could be. It shows the pattern isn't specific to MOs, not that it's neural.
\item One weak point against a probe cause: the fluctuations around the ramp aren't shared between areas.}

\q{Is it arousal, fatigue or satiety?}{
\item Beyond the shared trend, it doesn't track pupil, engagement, cumulative reward, reaction time or accuracy.
\item But ramp-shaped signals are hard to tell apart with these tests.}

\q{What about the outlier mouse, DY\_014?}{
\item It's dominated by one very high-rate cluster with one abrupt step: artifact-consistent.
\item Its effects go the same way, and dropping it changes nothing.}

\q{Did you change anything after seeing results?}{
\item One bug fix: a z-scoring floor that let a silent unit blow up held-out scores.
\item Raw primary $p$-values didn't move. Old outputs are kept.
\item Everything after the freeze is labelled exploratory or post hoc.}

\subsection*{Big picture}

\q{Isn't this just a failed hypothesis?}{
\item Yes, the central prediction failed. The proposal planned for that outcome.
\item What it taught: these methods rank by variance, a small belief signal loses to a large session ramp, and
we characterized that ramp.}

\q{Could the prior be coded dynamically rather than statically?}{
\item This design can't tell.
\item 30 s spans only 6--8 trials; blocks last about 46.
\item Both kinds of code are predictable 30 s ahead.
\item Analysis aligned to block switches would separate them.}

\q{What does this mean for predictive learning in ML?}{
\item Predictive objectives grab the biggest predictable signal first.
\item If a nuisance like drift is large, it takes the capacity, and small task variables can be missed.}

\q{How would you fix the method?}{
\item Add time or trend as a nuisance in the objective, use a drift-robust objective, or use more dimensions.
\item Test the fixes on a ground-truth simulation with a known drift and a known small prior.}

\q{What would change your conclusion?}{
\item A decoder or preprocessing where 30 s consistently beats 0.5 s for the prior across mice.
\item Matched-unit resampling showing the ramp is an artifact whose removal reveals a horizon-specific prior.}

\q{Did you use AI tools?}{
\item Yes, disclosed in the report as the course requires.}

\end{document}
